\chapter{Introduction} \label{cap:introduction}

The foreign exchange market (Forex or FX) trades roughly USD 7.5 trillion a day across spot, swap and derivative instruments \cite{bis2022triennial}. A market of that magnitude does not respond to price charts alone. Currency pairs react, often within seconds, to central-bank decisions, military escalation, sanctions, supply-chain shocks and political instability. Market participants able to ingest, interpret and act on such textual information faster than the average competitor retain a measurable advantage.

Trading strategies that rely exclusively on technical indicators applied to OHLC price bars \cite{murphy1999technical} have been documented to perform adequately under stable price regimes and to deteriorate when news flow, rather than price flow, drives the market. Modern Natural Language Processing (NLP) and the public release of finance-tuned language models \cite{araci2019finbert,yang2020finbert,wu2023bloomberggpt,liu2023fingpt} have substantially lowered the cost of revisiting that limitation empirically. The present work investigates a hybrid trading pipeline that complements price-based signals with sentiment- and event-based features extracted from text, evaluated under a strict backtesting protocol on historical data.

\section{Theme and Problem} \label{sec:theme-problem}

The theme is the integration of textual information (news articles, institutional reports, macroeconomic announcements) into algorithmic decision-making for Forex. The driving research question is the following:

\begin{quote}
Can NLP-derived textual features improve risk-adjusted Forex performance over a price-only baseline?
\end{quote}

\subsection*{Hypothesis} \label{sec:hypothesis}
The central premise of this study, that adding NLP-derived sentiment and event features to a competitive price-only baseline improves risk-adjusted Forex performance, is a \textit{hypothesis}, not an established result. It is rendered plausible by the equity-market evidence reviewed in Section~\ref{sec:justification} and by the documented failure of a macro-only fusion on precisely the two target pairs \cite{vanhelden2021fx}; plausibility, however, is not proof. The empirical apparatus developed in this work exists to subject that premise to a falsifiable test rather than to presuppose it. Accordingly, the research question above is operationalized as a single falsifiable hypothesis:
\begin{description}
  \item[\textbf{H1}] On EUR/USD and USD/JPY over the 2024 test span,
  the hybrid strategy of Chapter~\ref{cap:methodology} produces a higher
  annualized Sharpe ratio than the price-only baseline trained on the same
  data through the same backtesting protocol, by a margin $\Delta$ that
  survives the empirical transaction-cost model of Section~\ref{sec:evaluation}
  and a statistical test of return-distribution difference at the conventional
  significance level. Equivalently, the hybrid strategy delivers a smaller
  maximum drawdown under the same protocol. Either condition independently is
  sufficient to accept~H1.
\end{description}
H1 is referenced under an explicit accept/reject criterion in
Section~\ref{sec:final-remarks} of the conclusion, derived from the evidence
reported in Chapter~\ref{cap:experimental-evaluation}. The test span actually
delivered there is a reproduced one-year window (March 2016--February 2017),
not the 2024 span named above; Chapter~\ref{cap:experimental-evaluation}
states this substitution explicitly and reports its results as exploratory
rather than as a completed test of H1 in its originally scoped form.
A rejection constitutes
a valid and reportable outcome: it would establish that, under this protocol
and on these pairs, the textual channel does not yield a statistically
meaningful improvement over the price-only baseline. Such a negative finding
is itself a contribution, and motivates stating the hypothesis in a form that
the evidence can refute.

The problem decomposes into three sub-problems, each a distinct methodological challenge. First, a data problem: text and events at currency-grade scale must be collected from multiple sources, deduplicated, and time-aligned with FX prices \cite{leetaru2013gdelt}. Second, a representation problem: free-form text must be transformed into numeric features that retain financially meaningful signal (polarity, intensity, topical focus, uncertainty) without leaking future information into past decisions \cite{loughran2011liability,tetlock2007giving}. Third, a decision problem: the combined feature space must drive entries, exits and position sizing in a manner that remains auditable and robust to regime change \cite{lopezdeprado2018advances}.

\section{Justification and Motivation} \label{sec:justification}

Three lines of reasoning motivate this investigation.

The empirical case is decades old, although most of the supporting evidence comes from equities rather than Forex. \citeonline{tetlock2007giving}, working on the US equity market, demonstrated that media tone has predictive content beyond what is contained in past prices. \citeonline{caldara2022geopolitical} constructed a geopolitical risk index from newspaper text and documented measurable real-activity and broad financial-market effects following GPR shocks. Both contributions --- Tetlock on US equities and Caldara and Iacoviello on macroeconomic and broad-market outcomes --- exemplify the dominant lines of text-based financial prediction; Forex itself has been studied predominantly through macroeconomic and microstructure variables, with comparatively little use of textual inputs. This asymmetry leaves room for an incremental contribution and motivates this work as a transposition of an equity-derived hypothesis into a Forex setting.

Two recent USP monographs sharpen this gap on both sides. On the model side, \citeonline{felix2024swing}, produced in the same ICMC/USP MBA in Artificial Intelligence and Big Data as the present work, surveys seven machine-learning classifiers for swing-trading on the IBrX-100 and explicitly identifies textual information (``new financial articles, social media and even publications from economic analysts'') as a relevant qualitative input class that the implementation nevertheless leaves out. On the asset-class side, \citeonline{vanhelden2021fx} builds a hybrid FX strategy that fuses a macro-fundamental signal (real interest-rate differential, GDP-growth differential and M2/reserves ratio) with a technical signal averaged from three moving-average crossovers, applied to twelve currency pairs. The hybrid produces an annualized Sharpe ratio above 1 on all six emerging Latin-American pairs in the sample, but \textbf{drops to 0.24 on EUR and to $-1.08$ on JPY against the USD} (\citeauthor{vanhelden2021fx}, Table~6, p.~35). The two pairs on which the macro-only hybrid fails are precisely those targeted by the present work; the failure of the macro-only fusion on the developed-market pairs constitutes, in itself, a direct empirical motivation to replace the macro-fundamental channel with an NLP-derived textual channel.

The technological case is recent. Until 2018, financial text analytics relied predominantly on dictionary methods such as Loughran--McDonald \cite{loughran2011liability}, whose sensitivity is bounded by the contents of the lexicon. Transformer-based models \cite{vaswani2017attention,devlin2019bert} and their finance-tuned descendants \cite{araci2019finbert,yang2020finbert,wu2023bloomberggpt,liu2023fingpt} have advanced the achievable accuracy on financial sentiment well beyond the lexicon ceiling. Revisiting the Forex sentiment problem under this newer toolset is empirically tractable, and brings the resulting methodology closer to what is operationally feasible.

A third motivation is methodological reusability. Conducted as an applied dissertation within an MBA in Artificial Intelligence and Big Data, the present study is organized so that its methodology and pipeline remain extensible. Design decisions accordingly favor permissively-licensed components, an explicit interface boundary between perception, decision and execution, an auditable and configurable filter chain, and a perception layer decoupled from execution determinism so that model upgrades do not compromise reproducibility. The objective is a reusable methodology, not a one-off study: subsequent investigations should be able to refine the sentiment models, extend the asset universe or alter the evaluation protocol without reimplementing the apparatus.

\section{Objectives} \label{sec:objectives}

\subsection{General objective}

To design, implement and empirically evaluate a hybrid algorithmic-trading methodology for the Forex market, encompassing the data pipeline, the sentiment- and event-feature layers, the baseline and hybrid strategies and the backtesting protocol, applied to two major Forex pairs, EUR/USD and USD/JPY, and to determine, under controlled experimental conditions, whether the textual features yield a statistically meaningful improvement over a price-only baseline.

\subsection{Specific objectives}

\begin{enumerate}
    \item \textbf{Construct a reproducible data layer.} Assemble at least a decade of intraday price history for the selected pairs, complemented by a multi-source historical corpus of news and event data (GDELT \cite{leetaru2013gdelt} as the primary source, with additional news archives integrated through a uniform adapter interface) and macroeconomic-announcement calendars, under strict point-in-time discipline.

    \item \textbf{Define a sentiment- and event-feature layer.} Score textual and event information at the same temporal granularity as price data, employing FinBERT-class language models \cite{araci2019finbert,yang2020finbert} for sentence-level polarity, GDELT events for event-level features, and the Caldara--Iacoviello GPR index \cite{caldara2022geopolitical} as a continuous risk-regime signal.

    % NOTE (TODO/features): the specific feature and filter set is still under selection
    % in Chapter 3 (Methodology). The examples cited in this objective are representative,
    % not definitive; update both the lists below and the corresponding methodology section
    % once the filter chain is finalized. Keep this objective in sync with Chapter 3.
    \item \textbf{Specify and implement baseline and hybrid strategies in a shared framework.} The baseline draws on price-only feature families spanning technical indicators (such as moving averages, momentum oscillators and volatility measures), chart and candlestick patterns, and market-activity signals (such as volume and realized volatility). These are combined through per-family LightGBM sub-models whose calibrated probabilities feed a logistic meta-learner; the meta-learner output then drives a deterministic filter chain that emits BUY, SELL or HOLD on each bar. The hybrid retains the same architecture and adds a news- and event-feature family (FinBERT-class sentence-level scoring, GDELT events, and the Caldara--Iacoviello GPR index) as an additional sub-model in the meta-learner. Implementation proceeds in stages, from a minimum viable baseline up to the full architecture, with the exact family decomposition and filter set subject to revision based on experimental ablations (Section~\ref{subsubsec:phased-delivery}). Both strategies are evaluated within the same backtesting protocol, incorporating realistic transaction costs (spreads, slippage and rollover).

    \item \textbf{Conduct a rigorous empirical evaluation.} Apply out-of-sample risk-adjusted metrics (Sharpe, Sortino, maximum drawdown, profit factor), formal statistical tests of return-distribution differences, and ablation analyses isolating the marginal contribution of the sentiment and event channels.
\end{enumerate}

\section{Document Structure} \label{sec:structure}

The remainder of this dissertation is organized as follows. Chapter~\ref{cap:theoretical-foundation} reviews the theoretical foundations: the structure of the Forex market, the principles of algorithmic and quantitative trading, the NLP concepts on which the methodology depends, the literature on geopolitical risk and event data, and recent advances in finance-tuned language models. Chapter~\ref{cap:methodology} presents the methodology: data sources, feature engineering, model choices, backtesting protocol and evaluation criteria. Chapter~\ref{cap:experimental-evaluation} reports the experimental results, including the ablation analyses that isolate the contributions of the sentiment and event channels. Chapter~\ref{cap:conclusion} synthesizes the findings, discusses the limitations, and outlines directions for subsequent research.

The theoretical foundations and the methodology are presented in their definitive form. The experimental chapter reports the empirical work completed to date: the experimental setup, the baseline and hybrid strategy results, the confirmation reruns under the execution model, the ablation sweeps, the one-year and one-trading-year protocols including a fully reported news-only arm with registered paired inference, the adaptive-retraining frozen-policy result, the candlestick-catalog screen with its multiple-comparisons check, a reported data-coverage table with its three known and explicitly unresolved gaps (Section~\ref{sec:coverage-results}), and the comparison with the closest prior art. The confluence chain and the second exogenous (intermarket) feature family, two of the four follow-up studies registered alongside this work, remain fully outstanding at the time of submission; the price-action/momentum/Ichimoku contest, the support-and-resistance extension, and four of the five adaptive-retraining policy arms also remain future work. These are reported as such in Chapter~\ref{cap:conclusion} rather than filled in with unsupported numbers.
